\subsection{An explicit sufficient range}\label{sec:effective}
The sign estimate alone does not give a cutoff for Proposition~\ref{prop:index}: the collision and deletion costs must also be small compared with $\gamma$. Here are numerical upper bounds from the preceding counts. They are deliberately loose. Write $L=L_p$, $J=\min(I,\sqrt p)$. On the stated progression we will show that
\begin{equation}\label{eq:explicit-index}
 \mathcopy{log(p) >= 2^256 and I >= L_p^24 imply m >= 2^(-135) min(I, sqrt(p)) / (L_p^3).}{\log p\ge 2^{256},\quad I\ge L^{24}
 \quad\Longrightarrow\quad m\ge \frac{J}{2^{135}L^3}}
\end{equation}
for the number $m$ of added opposite pairs. There are at most $4L$ dyadic direction bins. In each bin there are at most $8D^2$ directions and $4T^2$ ordered parameter pairs; $DT\le4p$. Thus across all bins at most $2^{11}p^2L$ ordered templates enter any one-root count.

\paragraph{One-index costs.} The anchor tests in Section~\ref{sec:sampling} cost at most $(2\binom82+8)/p=64/p$. A same-chord collision costs at most $2^{11}hL/p$. For Lemma~\ref{lem:lifts}, fix $E,d$ and one bin. If $D\le U$, the progression bound gives at most $8D^2(1+32U/D)\le264DU$ solutions; if $U<D$, the bound is $\sum_{m\le U}8m(1+64D/m)\le520DU$, with $m\ge1$. Both are at most $2^{12}p$ since $DU=4p/d$. Allow fewer than $2^{17}$ matrices, at most $2^7$ choices of $d$, $4L$ bins and a factor four for column orders. Dividing by $p^2$ gives at most $2^{40}L/p$.

For Lemma~\ref{lem:gcd}, use $1+\log(4p)\le3L$. The squared residue weights sum to at most $18(4p)(3L)^3<2^{11}pL^3$, and either signed small-factor sum is at most $24pL/\ell$. Including the factor 16 for signed splittings and all four rows bounds the large-gcd probability by $2^{30}L^5/G$, even allowing four endpoint designations. For the small gcds in Proposition~\ref{prop:nonparallel}, use its variables $X,Y$. When $X\ge p$, the count per bin pair is at most
\[
  16Y^2\cdot 2h\,[128(1+G)X/p]^2.
\]
Here the rectangle has area parameter $4X$, and the factor 2 is the quadratic root count. For $X<p$ the analogous bound is $64X^2 h[32(1+G)Y/p]^2$, using $Y\ge p$. With $XY\le16p^2$, $(4L)^2$ bin pairs and four designations, the small-gcd probability is at most $2^{34}(h/p)(1+G)^2 L^2$. Now $G=\lfloor I^{1/4}\rfloor\ge I^{1/4}/2$, $1+G\le2I^{1/4}$ and $p\ge I\ge L^{24}$. The sum of all the one-index costs just listed is at most $2^{50}/L$.

\paragraph{Two- and three-index costs.} For one choice of signs at two indices, the mean of the off-diagonal inverse direction height is at most $8(4p+p(1+\log(4p)))/p^2\le56L/p$: at height $D$ there are at most $8D$ directions, each with at most $1+p/D$ positions in the second sampling square. The carrier cost is at most $1/p$ plus four times this mean. For all four signs, including shared coordinates, the eight anchors and the mate diameters, the geometric cost is at most $2^{11}L/p$ (each fixed line or coordinate costs at most $1/p$). In Proposition~\ref{prop:cross} the nonexceptional cost is at most $4(2^{11}p^2L)^2h/p^5$. In the exceptional case the two bounds per second bin are $288D_1^2T_1$ and $136D_1T_1^2$: the respective multiplicities are $8\cdot2\cdot2\cdot9$ and $4\cdot2\cdot17$. Their minimum is at most $2^{12}p^{3/2}$. Multiply by $2^{11}p^2L$ first templates, $4L$ second bins, four designations and $p$ possible outcomes, then divide by $p^5$. Together the two-index costs are at most $2^{31}L^2/J$. The three-index cost is at most $8/p$.

Finally, $N\le J/L^3$ gives $\mathbb E Z/N\le2^{31}/L+8/L^6$. The size assumption in~\eqref{eq:explicit-index} more than meets the sign-test threshold of Lemma~\ref{lem:signs}, so we may use $\gamma=2^{-133}$. Its role here is to make the collision losses small compared with $\gamma$. The combined loss per proposal is at most $2^{51}/L<\gamma/2$. Since $J\ge L^{12}$, taking the floor gives $N\ge J/(2L^3)$. The same deletion and matching as before therefore retain at least $\gamma N/2\ge J/(2^{135}L^3)$ pairs. This proves~\eqref{eq:explicit-index}; it is a sufficient range, not a prediction of the first successful prime.
